% \section{Dataset: HistoriQA-ThirdRepublic}
% \label{sec:dataset}

% We build on HistoriQA-ThirdRepublic~\cite{historiqa}, a French
% historian-validated benchmark over three heterogeneous sources from
% 1887: the parliamentary transcripts of the Chambre des Députés
% (\textit{Les Débats}) and two ideologically opposed dailies
% (\textit{Le Gaulois}, \textit{L'Intransigeant}). The benchmark contains
% 1{,}782 questions, of which 885 are multi-hop; the latter split into
% 571 newspaper~$\to$~\textit{Les Débats} pairs and 314 cross-newspaper
% pairs. Full corpus statistics, historical context, and question-type
% breakdown are deferred to Appendix~\ref{sec:appendix-dataset}
% (Table~\ref{tab:corpus-stats}).

% \paragraph{Retrieval challenges that shape our method.}
% Three structural properties directly motivate the task-conditioned
% framework introduced in Section~\ref{sec:methodology}:
% \begin{itemize}[leftmargin=1.2em,itemsep=0.2em,topsep=0.3em]
%   \item \emph{Length heterogeneity and OCR noise.} Parliamentary
%     transcripts average an order of magnitude more tokens per document
%     than newspaper articles, so no single chunking granularity fits all
%     sources.
%   \item \emph{Corpus imbalance.} \textit{Les Débats} is more than six
%     times larger than the two newspapers combined and dominates a
%     unified dense retriever, supplying 80.6\% of top-3 hits in the naive
%     baseline~\cite{historiqa} even when the gold evidence lies entirely
%     in the press.
%   \item \emph{Cross-source multi-hop reasoning.} 314 questions require
%     both gold passages from the newspapers and 571 require at least one
%     newspaper and one parliamentary passage  - a structure that plain
%     top-$k$ cosine ranking cannot guarantee, since the dominant corpus
%     may consume all slots.
% \end{itemize}
% The first challenge motivates per-cluster chunking
% (Sec.~\ref{sec:indexing-results}); the latter two motivate the source
% routing (QRe) and uniform multi-source (UMS) strategies of
% Section~\ref{sec:retrieval-results}.


\section{Dataset: HistoriQA-ThirdRepublic}
\label{sec:dataset}

We build on HistoriQA-ThirdRepublic~\cite{historiqa-pellet}, a French
historian-validated benchmark over three heterogeneous 1887 sources
digitised by the Biblioth\`eque nationale de France and distributed
through Gallica\footnote{\url{https://gallica.bnf.fr}}:
the parliamentary transcripts of the Chambre des D\'eput\'es
(\textit{Les D\'ebats}) and two ideologically opposed dailies
(\textit{Le Gaulois}, \textit{L'Intransigeant}). The benchmark
contains 885 historian-validated multi-hop questions, split into
571 newspaper-\textit{D\'ebats} pairs and 314 cross-newspaper
pairs. Full corpus statistics, historical context, and the multi-hop
question construction pipeline are in Appendix~\ref{sec:appendix-dataset}
(Table~\ref{tab:corpus-stats}).

\paragraph{Retrieval challenges that shape our method.}
Two structural properties drive the framework: (i)~\emph{corpus
heterogeneity and imbalance}: \textit{Les D\'ebats} is more than
$30\times$ larger than the two newspapers combined and its individual
documents are an order of magnitude longer than newspaper articles,
which both rules out a single chunking granularity (motivating
per-cluster chunking, Section~\ref{sec:indexing-results}) and lets
\textit{D\'ebats} supply 80.6\% of Naive-baseline top-3 hits even on
newspaper-only gold; and (ii)~\emph{cross-source multi-hop reasoning}:
the multi-hop questions above require evidence from multiple corpora,
a structure that top-$k$ cosine ranking cannot guarantee since the
dominant corpus may consume all slots. Together, these properties
motivate the source-routing (QRe) and uniform multi-source (UMS)
strategies of Section~\ref{sec:retrieval-results}.